% !TEX root = ../journal_v2_0421052099.tex
%
% Condensed from thesis_book/chapters/environment.tex and the knowledge-graph
% preliminaries of thesis_book/chapters/background.tex.
\section{The Knowledge Environment}\label{sec:environment}

Everything the agent can know, it knows from the graph. So the environment sets
a ceiling on how accurate any system can be. A result reported without that
ceiling is hard to interpret.

\subsection{Freebase and Mediator Nodes}

A knowledge graph $G = (E, R, T)$ stores facts as triples $(h, r, t)$, each
asserting that relation $r$ holds from head entity $h$ to tail entity $t$.
Triples are atomic, so a claim can be checked by testing for specific edges.
They are also composable, and a chain of them expresses a multi-hop relationship
that no single triple states. Freebase~\cite{bollacker2008freebase}, frozen
since 2016, is the substrate of the standard benchmarks. One of its features
changes every depth measurement. A \emph{mediator node} is an unnamed node
standing for one fact with more than two participants, with one edge to each
participant. That Rainn Wilson played Dwight Schrute in \emph{The Office} is
stored as a mediator with edges to the actor, the character, and the show. So a
relationship a user would describe as one hop is two edges in the graph.
Distances in this paper are in \emph{raw} hops, which count graph edges. One
\emph{logical} hop through a mediator is two raw hops.

\subsection{Hallucination, Defined Operationally}\label{sec:halluc}

``Hallucination'' means different things in different papers. This paper adopts
a narrow, operational definition tied to the environment. An \emph{asserted}
entity is one the system names in its answer. It is \emph{ungrounded} if it does
not exist in the graph, or exists but lies more than four hops from every topic
entity of the question. A \emph{claim} is a single statement drawn from the
answer that relates two entities. It is \emph{unsupported} if the traversed
triples do not entail it. Two consequences follow. The definition is relative to
the environment, so an entity absent from the graph is ungrounded even if it is
true in the world. And groundedness is not correctness. An answer can be
grounded in every entity and still be the wrong answer to the question.

\subsection{Construction}

The environment is the union of every per-question subgraph that the
Reasoning-on-Graphs distribution~\cite{luo2024reasoning} publishes for WebQSP
and ComplexWebQuestions, across all splits. So the questions and the graph come
from one distribution and are mutually consistent. Taking the union is not a
minor detail. A per-question subgraph on its own is close to an oracle. It was
extracted to contain the answer, so an agent searching it meets almost no
distractors. Merging every question's neighbourhood into one graph brings the
distractors back. \Cref{tab:environment} records the result. Mediators are
retained rather than collapsed, since systems that skip them lose the facts they
encode~\cite{xu2024generate}. They make up $63.7\%$ of the nodes.

Two consequences must be disclosed. The subgraphs were pre-extracted to
guarantee that answers are reachable, which makes the environment friendlier
than raw Freebase. And the distribution keys entities by name rather than by
Freebase identifier. So distinct entities that share a surface form are merged
into one node. Both bound the absolute numbers, and \Cref{sec:threats} returns
to them.

The graph is stored in Neo4j~\cite{robinson2015graph}, with each relation as a
native relationship type, and is queried in Cypher~\cite{francis2018cypher}.
Entity names and the relation vocabulary are embedded with the 384-dimensional
sentence encoder all-MiniLM-L6-v2~\cite{reimers2019sentence}. Relation names are
verbalised before embedding, so \textsf{people.person.place\_of\_birth} is
encoded as the phrase it means. Embeddings have a narrow role. They link
question phrases to entities and pre-rank relation candidates. They never
establish that a fact is true, a role that belongs to the triples alone.

Entity resolution is a three-stage cascade of exact, lexical, and vector
matching, each stage tried only if the previous one fails. Over all $7{,}149$
topic-entity mentions in both test sets, $99.86\%$ resolve at the exact stage.
The ten that do not are bare machine identifiers that the benchmark carries in
place of a name. So a failure in this environment cannot be blamed on the
resolver missing the question's topic. The graph and the benchmarks share their
name strings, though, so this shows only that the resolver loses nothing the
environment contains.

\begin{table}[!t]
\caption{The knowledge environment and its validation gate. Reachability is the
  share of test questions with a gold answer within four raw hops of a topic
  entity. It is the Hits@1 ceiling over each full test split.}
\label{tab:environment}
\centering
\small
\begin{tabular}{@{}lr@{}}
\toprule
Property & Value \\
\midrule
Entities & $2{,}592{,}892$ \\
Triples & $8{,}309{,}194$ \\
Distinct relation types & $7{,}058$ \\
Mediator nodes & $1{,}652{,}618$ ($63.7\%$) \\
Mentions resolved by exact match & $7{,}139$ of $7{,}149$ ($99.86\%$) \\
\midrule
Reachable, WebQSP test split & $1{,}584$ of $1{,}628$ ($97.3\%$) \\
Reachable, CWQ test split & $3{,}519$ of $3{,}531$ ($99.7\%$) \\
\bottomrule
\end{tabular}
\end{table}

\subsection{The Validation Gate}\label{sec:gate}

No experiment ran until the environment was certified. For every test question
we computed whether a gold answer is reachable from a topic entity within four
raw hops, by breadth-first search over undirected edges. The agent's tools look
in both directions, so the gate does the same. Reachability is $97.3\%$ on
WebQSP and $99.7\%$ on CWQ (\Cref{tab:environment}). It is the accuracy ceiling,
since no navigating system can find an answer it cannot reach. Existence and
reachability agree question for question. When a question is unanswerable here,
it is because its gold entities are absent, not because they are too far away.
Reachability was also re-verified directly in Cypher, independently of the code
that computed it.

Certifying reachability says nothing about answers that are not entities. The
import gives every node a single label, with no typed literal. Only $77$ nodes
carry a full date as their name and $175$ a bare year, together under $0.01\%$
of the graph and untyped. And the tool layer exposes no sorting or comparison
operator. So ``the earliest'' or ``the smallest'' cannot be computed even when
every candidate is present. These gaps define a class of questions that are
\emph{unanswerable in this environment}, those whose answer is a date, a
measured quantity, or a superlative selection. Missing one is not a reasoning
failure, and \Cref{sec:errors} counts how often it happens.

\endinput
